\documentclass[../../../include/open-logic-section]{subfiles}

\begin{document}

\olfileid[id]{sth}{story}{approach}
\olsection{Pendekatan Kumulatif-Iteratif}

Berikut pernyataan yang sedikit lebih lengkap tentang cara kita akan
menstratifikasi himpunan ke dalam tahap-tahap:
\begin{quote}
	Himpunan dibentuk dalam \emph{tahap-tahap}. Untuk setiap tahap $S$,
	terdapat tahap-tahap tertentu yang berada \emph{sebelum} $S$. Pada tahap
	$S$, setiap koleksi yang terdiri atas himpunan yang dibentuk pada tahap
	sebelum $S$ dibentuk menjadi sebuah himpunan. Tidak ada himpunan selain
	himpunan yang dibentuk pada tahap-tahap tersebut.
	\citep[p.~323]{Shoenfield:AST}
\end{quote}
Ini merupakan sketsa \emph{konsepsi kumulatif-iteratif tentang himpunan}.
Konsepsi ini akan mendasari teori himpunan formal yang kita sajikan dalam
\olref[sth][][]{part}.

Mari kita telaah hal ini dengan sedikit lebih terperinci. Sebagaimana
Shoenfield menggambarkan prosesnya, pada setiap tahap kita membentuk himpunan
baru dari {himpunan} yang telah tersedia dari tahap-tahap sebelumnya. Jadi,
dalam gambaran Shoenfield, pada tahap awal, yaitu tahap $0$, tidak ada tahap
yang \emph{lebih awal}; karena itu, \emph{a fortiori}, tidak ada himpunan yang
tersedia bagi kita dari tahap sebelumnya.\footnote{Mengapa kita harus
mengasumsikan bahwa \emph{ada} tahap pertama? Lihat catatan kaki untuk
\stagesord{} dalam \olref[z][story]{sec}.} Jadi, kita hanya membentuk satu
himpunan: himpunan tanpa !!{element}s, $\emptyset$. Pada tahap $1$, tepat satu
himpunan tersedia dari tahap sebelumnya, sehingga hanya ada satu himpunan
baru, yaitu $\{\emptyset\}$. Pada tahap $2$, dua himpunan tersedia dari tahap
sebelumnya, dan kita membentuk dua himpunan baru,
$\{\{\emptyset\}\}$ dan $\{\emptyset, \{\emptyset\}\}$. Pada tahap $3$,
empat himpunan tersedia dari tahap sebelumnya, sehingga kita membentuk dua
belas himpunan baru\ldots. Dengan demikian, gambaran kumulatif-iteratif
tentang himpunan akan tampak kira-kira seperti ini (angka menunjukkan tahap):
\begin{center}
	\begin{tikzpicture}[scale=0.6]
	\tikzset{cut_here/.style={densely dotted}}
	\draw (-4,6) -- (-1, 0)-- (2,6);
	\draw[cut_here] (-5,8)--(-4,6);
	\draw[cut_here] (2,6)--(3,8);
	\draw(-1.5, 1)--(-0.5, 1);
	\draw(-2, 2)--(0, 2);
	\draw(-2.5, 3)--(0.5,3);
	\draw(-3, 4)--(1, 4);
	\draw(-3.5, 5)--(1.5, 5);
	\draw(-4, 6)--(2, 6);
	\node[label] at (0, 0) {\small 0};
	\node[label] at (0.5, 1) {\small 1};
	\node[label] at (1, 2) {\small 2};
	\node[label] at (1.5, 3) {\small 3};
	\node[label] at (2, 4) {\small 4};
	\node[label] at (2.5, 5) {\small 5};
	\node[label] at (3, 6) {\small 6};
	\end{tikzpicture}
\end{center}
Jadi, mengapa kita perlu menerima cerita ini?

Salah satu alasannya ialah bahwa cerita ini baik dan mudah dikelola. Setelah
cerita yang paling jelas, yaitu Komprehensi Naif, tumbang, kita memerlukan
sesuatu yang baik.

Namun, cerita ini bukan \emph{sekadar} baik. Kita memiliki alasan kuat untuk
percaya bahwa setiap teori himpunan yang didasarkan pada cerita ini akan
\emph{konsisten}. Berikut alasannya.

Menurut konsepsi kumulatif-iteratif tentang himpunan, kita membentuk himpunan
pada tahap-tahap tertentu; dan !!{element}s-nya harus berupa objek yang
\emph{sudah} tersedia. Jadi, untuk setiap tahap~$S$, kita dapat membentuk
himpunan
\[
	R_S = \Setabs{x}{x \notin x \text{ dan $x$ tersedia sebelum $S$}}
\]
Penalaran dalam pembuktian Paradoks Russell kini akan menetapkan bahwa $R_S$
sendiri tidak tersedia sebelum tahap $S$. Dan itu bukan kontradiksi. Selain
itu, jika kita menerima konsepsi kumulatif-iteratif tentang himpunan, kita
bahkan tidak semestinya \emph{mengharapkan} dapat membentuk himpunan Russell
itu sendiri. Sebab, himpunan itu akan menjadi himpunan semua himpunan yang
bukan anggota dirinya sendiri dan yang ``kelak akan tersedia''. Singkatnya:
fakta bahwa kita (secara terbukti) tidak dapat membentuk himpunan Russell
bukanlah sesuatu yang \emph{mengejutkan} menurut cerita kumulatif-iteratif;
itulah yang akan kita \emph{prediksi}.

%In one sense, then, the cumulative-iterative conception of set yields a response to Russell's Paradox which is rather like the \emph{predicativist}'s response. After all, the predicativist said that the collection of all non-self-membered sets$_n$ is a set$_{n+1}$, and this set$_n$ will not be self-membered. (Indeed, most predicativists will treat it as \emph{ungrammatical} to try to ask whether a set is self-membered.)
%
%But there is an important difference between the cumulative-iterative approach and the predicativist approach: the cumulative-iterative approach treats all of our entities as being \emph{of the same kind}. The predicativist will presumably have an empty set$_0$ (i.e., a set$_0$ with no !!{element}s), and an empty set$_1$ (i.e., a set$_1$ with no !!{element}s), and these will be \emph{different entities}. But on the cumulative-iterative approach, there is just one empty set, $\emptyset$, and it will be available at every stage of the hierarchy.

%Indeed, the Axiom of Extensionality, stated at the start of this chapter, will hold true of the elements in this hierarchy (so that there can be \emph{only one} ``empty'' set). But I do not intend to use the cumulative-iterative conception to justify Extensionality (see \cite{Boolos1971}). Again: I suggest that we should accept Extensionality, just because we are interested in extensional collections.

\end{document}
